\chapter{Numerical Integration}

In the previous chapter, we applied integration across geometry and physics:
volumes of solids, arc lengths, work, and probability distributions. In each
case, the application led to a definite integral $\int_a^b f(x)\,\dd x$.

We usually expect to evaluate such integrals using the Fundamental Theorem of
Calculus by finding an antiderivative $F$ with $F' = f$. But that requires being
able to express $F$ in closed form.

For many functions, no such antiderivative exists. For example,
\[
\int_0^1 e^{-x^2}\,\dd x
\]
has no antiderivative in terms of elementary functions---no finite combination
of polynomials, roots, exponentials, logarithms, and trigonometric functions.
Joseph Liouville proved in the 1830s that an elementary antiderivative for
$e^{-x^2}$ does not exist. Yet the integral represents a well-defined area
(approximately $0.7468$), and it appears throughout probability and statistics
as the normal distribution.

When an antiderivative cannot be found, we return to the definition of the
integral as an accumulated sum. By evaluating the integrand at selected sample
points, we can estimate the area to any desired accuracy with certified error
bounds. That is the goal of numerical integration.

\begin{focusbox}
The integral $\int_a^b f\,\dd x$ is a number: the signed area under $f$. The Fundamental Theorem is one road to that number, available only when an antiderivative exists in usable form. Numerical integration is a different road, always open, that returns to the original meaning of the integral as accumulated area and computes that area directly from samples of $f$. Our goal: get an accurate number \emph{and} a reliable bound on how wrong it might be.
\end{focusbox}

\section{When the Antiderivative Will Not Come}
\label{ch27:noanti}

Numerical methods are necessary for two distinct reasons: functions with no
elementary antiderivative, and functions known only through discrete data points.

Here are integrals that arise constantly and have no elementary antiderivative:
\[
\int e^{-x^2}\,\dd x,\qquad
\int \frac{\sin x}{x}\,\dd x,\qquad
\int \frac{1}{\ln x}\,\dd x,\qquad
\int \sqrt{1 + \tfrac14\cos^2 x}\,\dd x,\qquad
\int \frac{e^{x}}{x}\,\dd x.
\]
The first is the Gaussian in probability. The second, $\sin(x)/x$, arises in optics
and signal processing. The third, $1/\ln x$, relates to the distribution of prime numbers.
The fourth gives the arc length of an ellipse. None of them yields to substitution,
parts, or partial fractions, because no elementary antiderivative exists to be found.

\begin{whybox}
\textbf{Why an impossible antiderivative is not your fault.}
Joseph Liouville (1833--1841) proved which functions can have elementary
antiderivatives, establishing that $e^{-x^2}$, $\sin(x)/x$, and many others
provably cannot. The absence of an elementary antiderivative is an intrinsic
property of these functions, not a lack of technique. In such cases, numerical
integration is the standard path to a quantitative answer.
\end{whybox}

And even when an antiderivative \emph{does} exist, you may not have a formula for $f$ at all. A bridge engineer measuring the deflection of a beam, a physiologist tracking cardiac output, a data scientist with a column of numbers -- they have $f$ only as a \emph{table of values}, sampled at points. There is no symbolic $f$ to integrate. There is only data, and the area under it. Numerical integration is the language for turning samples into accumulated quantities.

\begin{applicationbox}
A car's speedometer reads its speed $v(t)$ each second. How far did the car travel between $t = 0$ and $t = 60$? Distance is $\int_0^{60} v(t)\,\dd t$. You have no formula for $v(t)$ -- you have sixty-one numbers. The only way to recover distance is to estimate the area under the speed curve from those samples. Every fitness tracker, every flight data recorder, every digital oscilloscope integrates this way thousands of times a second.
\end{applicationbox}

\section{The Crudest Defensible Estimate: Rectangles}
\label{ch27:rectangles}

We already met the idea in the chapter on the integral. Chop $[a,b]$ into $n$ equal slices of width
\[
h = \frac{b-a}{n},
\]
with cut points $x_0 = a,\ x_1 = a+h,\ \dots,\ x_n = b$. On each slice approximate the area by a rectangle. If we use the left endpoint's height $f(x_{i-1})$ on each slice we get the \emph{left Riemann sum}; the right endpoints give the right sum.

The trouble with rectangles is visible to the eye. A rectangle has a flat top, but $f$ usually slopes. On a rising function, left rectangles sit entirely below the curve and lose area; right rectangles poke above and gain it (Figure~\ref{ch27:fig-rect}). The error on each slice is roughly the little triangular sliver between the flat top and the sloping curve. Halve the slice width and you halve roughly each sliver's base \emph{and} its height -- but there are twice as many of them, so the total error only halves. This is the painful arithmetic of rectangles: error $\propto h$. To get one more correct decimal digit you need ten times as many slices.

\begin{figure}[H]
\centering
\begin{tikzpicture}
\begin{axis}[
  width=11.5cm, height=6cm,
  domain=0:4, samples=100,
  axis lines=middle,
  xtick={0,1,2,3,4}, ytick=\empty,
  xlabel={$x$}, ylabel={$f(x)$},
  xmin=-0.3, xmax=4.5, ymin=0, ymax=3.2,
  enlargelimits=false,
]
\foreach \i in {0,1,2,3} {
  \edef\temp{\noexpand\addplot[draw=trapred!70, fill=trapred!12]
    coordinates {(\i,0) (\i+1,0) (\i+1,{0.6+0.25*\i+0.12*\i*\i}) (\i,{0.6+0.25*\i+0.12*\i*\i})} --cycle;}
  \temp
}
\addplot[curveblue, very thick, domain=0:4] {0.6 + 0.25*x + 0.12*x^2};
\node[curveblue, font=\small, anchor=south east] at (axis cs:2.9, 2.45) {$f(x)$};
\end{axis}
\end{tikzpicture}
\caption{Left rectangles under a rising curve. Each rectangle's flat top falls short of the sloping curve, leaving a triangular sliver of lost area. The total error shrinks only in proportion to the slice width $h$. We can do much better by letting the tops slope.}
\label{ch27:fig-rect}
\end{figure}

\begin{intuitionbox}
The guiding principle of numerical quadrature is to improve the local approximation of the curve. A rectangle's top is a constant (degree 0). A line segment (degree 1) accounts for the local slope, yielding the trapezoidal rule. A parabola (degree 2) accounts for local curvature as well, yielding Simpson's rule. Each higher-degree polynomial interpolation eliminates a leading error term, increasing the order of convergence.
\end{intuitionbox}

\section{The Trapezoidal Rule: Let the Tops Slope}
\label{ch27:trap}

The natural remedy for the rectangle error becomes clear once you inspect the missing sliver. Instead of a flat top, connect the two endpoints of each slice with a straight line. The flat-topped rectangle becomes a \emph{trapezoid} whose slanted top hugs the curve far better (Figure~\ref{ch27:fig-trap}).

\begin{figure}[H]
\centering
\begin{tikzpicture}
\begin{axis}[
  width=11.5cm, height=6cm,
  domain=0:4, samples=100,
  axis lines=middle,
  xtick={0,1,2,3,4}, ytick=\empty,
  xlabel={$x$}, ylabel={$f$},
  xmin=-0.3, xmax=4.5, ymin=0, ymax=3.2,
  enlargelimits=false,
]
\foreach \i in {0,1,2,3} {
  \edef\temp{\noexpand\addplot[draw=targetgreen!70, fill=targetgreen!12]
    coordinates {(\i,0) (\i+1,0) (\i+1,{0.6+0.25*(\i+1)+0.12*(\i+1)*(\i+1)}) (\i,{0.6+0.25*\i+0.12*\i*\i})} --cycle;}
  \temp
}
\addplot[curveblue, very thick, domain=0:4] {0.6 + 0.25*x + 0.12*x^2};
\end{axis}
\end{tikzpicture}
\caption{The trapezoidal rule: connect consecutive sample points with straight chords. Each trapezoid hugs the curve much more closely than a flat rectangle. The leftover error per slice now comes only from the curve's \emph{bending}, not its sloping.}
\label{ch27:fig-trap}
\end{figure}

The area of a single trapezoid is its average height times its width. On the slice $[x_{i-1}, x_i]$ the two heights are $f(x_{i-1})$ and $f(x_i)$, and the width is $h$, so
\[
\text{slice area} \approx h\cdot \frac{f(x_{i-1}) + f(x_i)}{2}.
\]
Add up all $n$ slices. Watch what happens: every interior point $x_1, \dots, x_{n-1}$ is the right endpoint of one trapezoid and the left endpoint of the next, so its height gets counted twice. The two outer endpoints $x_0$ and $x_n$ are counted once. The total telescopes into a clean formula.

\begin{focusbox}
\textbf{Trapezoidal Rule.} With $h = (b-a)/n$ and $x_i = a + i h$,
\[
T_n = \frac{h}{2}\Big[\, f(x_0) + 2f(x_1) + 2f(x_2) + \cdots + 2f(x_{n-1}) + f(x_n) \,\Big].
\]
The pattern of weights is $1,2,2,\dots,2,1$. Ends count once, interiors count twice, and the whole thing is scaled by $h/2$.
\end{focusbox}

\begin{whybox}
Why is the weight pattern $1,2,2,\dots,2,1$ and not something fancier? Because each interior sample is shared between two neighboring trapezoids -- it is the right wall of one and the left wall of the next -- so its height contributes to two pieces of area, hence weight $2$. The two boundary samples belong to only one trapezoid each, hence weight $1$. The trapezoidal rule is nothing more than ``replace the curve with the connect-the-dots polygon through your samples, and take the area under the polygon.'' That polygon is the unique piecewise-linear function agreeing with $f$ at every node (the sample points $x_i$ where we evaluate the function). Integrating it exactly gives $T_n$. So the trapezoidal rule is exact for any function that already \emph{is} a straight line on each slice -- and approximately right to the extent that $f$ is nearly straight at the scale of $h$.
\end{whybox}

\begin{recoverybox}
\textbf{Pause and think:} Suppose you use the trapezoidal rule to integrate $f(x) = 5x - 3$ from $x = 0$ to $x = 10$, using just a single slice ($n = 1$). What is the error?

\emph{Mental check:} The function is a straight line, so its second derivative is zero everywhere ($f'' = 0$). The top edge of the trapezoid is not an approximation of the line -- it \emph{is} the line. The error is identically zero, even with just one slice.
\end{recoverybox}

\begin{labbox}
Estimate $\int_0^1 e^{-x^2}\,\dd x$ with the trapezoidal rule using $n=4$. Then $h = 0.25$ and the nodes are $0, 0.25, 0.5, 0.75, 1$. Compute the heights $f(x) = e^{-x^2}$:
\[
\begin{array}{c|c}
x & e^{-x^2}\\\hline
0    & 1.000000\\
0.25 & 0.939413\\
0.5  & 0.778801\\
0.75 & 0.569783\\
1.0  & 0.367879
\end{array}
\]
Now apply the weights $1,2,2,2,1$:
\[
T_4 = \frac{0.25}{2}\Big[1.000000 + 2(0.939413) + 2(0.778801) + 2(0.569783) + 0.367879\Big].
\]
The bracket is $1.000000 + 1.878826 + 1.557602 + 1.139566 + 0.367879 = 5.943873$. Multiply by $h/2 = 0.125$: $T_4 = 0.742984$. The true value is $0.746824$. We are off by about $0.0038$ -- already under one percent with only five function evaluations, and we never found an antiderivative because none exists.
\end{labbox}

\subsection{How Wrong Is the Trapezoidal Rule? The Error Comes From Curvature}
\label{ch27:traperror}

Look again at Figure~\ref{ch27:fig-trap}. Where the curve bulges \emph{up} (concave down), the chord cuts underneath and the trapezoid undershoots. Where the curve bulges \emph{down} (concave up), the chord rides above and the trapezoid overshoots. Either way, the error is governed by how much the curve bends -- by the second derivative $f''$. A straight line has $f'' = 0$ and the trapezoidal rule nails it exactly. The more $f$ curves, the more the chord misses.

Make this precise on one slice, because the mechanism is the whole story. Take a single interval of width $h$ and ask how far the area under the chord differs from the area under $f$. A careful accounting (a short Taylor expansion -- approximating $f$ near the midpoint by its local value, slope, and bending -- done in the box below) gives, for one slice,
\[
\text{(true area)} - \text{(trapezoid area)} = -\frac{h^3}{12}\, f''(\xi)
\]
for some point $\xi$ inside the slice. The error on one slice is proportional to $h^3$ and to the curvature there.

\begin{happeningbox}
\textbf{Why $h^3$ per slice but $h^2$ overall.} Each slice carries error of size $\sim h^3$. But to cover a fixed interval $[a,b]$ you need $n = (b-a)/h$ slices, so the errors stack up:
\[
\text{total error} \sim n \cdot h^3 = \frac{b-a}{h}\cdot h^3 = (b-a)\,h^2.
\]
One power of $h$ is eaten by needing more slices as $h$ shrinks. What survives is $h^2$. This is the deep pattern of the whole subject: a method with per-slice error $h^{p+1}$ delivers \emph{global} error $h^{p}$. The trapezoidal rule has $p=2$. Halve $h$ (double the work) and the error drops by a factor of $4$, not $2$. That is already a world better than rectangles, where halving $h$ only halves the error.
\end{happeningbox}

\begin{whybox}
\textbf{Where the $-h^3 f''/12$ comes from.} Put the slice symmetrically as $[-h/2, h/2]$ to keep the algebra clean. Taylor-expand $f$ about the midpoint:
\[
f(x) = f(0) + f'(0)\,x + \tfrac12 f''(0)\,x^2 + \cdots
\]
Integrate across the slice. The odd term $f'(0)x$ integrates to zero by symmetry, leaving
\[
\int_{-h/2}^{h/2} f\,\dd x = f(0)\,h + \tfrac12 f''(0)\cdot \frac{2}{3}\Big(\tfrac{h}{2}\Big)^3 + \cdots = f(0)\,h + \frac{f''(0)\,h^3}{24} + \cdots
\]
The trapezoid uses the two endpoint heights $f(\pm h/2)$. Their average is
\[
\frac{f(h/2)+f(-h/2)}{2} = f(0) + \tfrac12 f''(0)\Big(\tfrac{h}{2}\Big)^2 + \cdots = f(0) + \frac{f''(0)h^2}{8} + \cdots,
\]
so the trapezoid area is $h\big(f(0) + \tfrac{f''(0)h^2}{8}\big) = f(0)h + \frac{f''(0)h^3}{8}$. Subtract:
\[
\text{true} - \text{trap} = \frac{f''(0)h^3}{24} - \frac{f''(0)h^3}{8} = f''(0)h^3\Big(\tfrac{1}{24} - \tfrac{3}{24}\Big) = -\frac{f''(0)h^3}{12}.
\]
There it is. The midpoint rule (rectangle of height $f(0)$, area $f(0)\,h$) misses the truth by $+\tfrac{1}{24}f''(0)h^3$, while the trapezoid misses by $-\tfrac{1}{12}f''(0)h^3$. Both errors are governed by $f'' h^3$, but the trapezoid's error is exactly
\emph{twice as large, in the opposite direction}. This relationship is the foundation
for Simpson's rule.
\end{whybox}

Summing $-h^3 f''(\xi_i)/12$ over all $n$ slices and replacing the average of the $f''(\xi_i)$ by a single value $f''(\xi)$ (using the Intermediate Value Theorem, which guarantees that a continuous function on an interval attains its average value) gives the standard bound.

\begin{focusbox}
\textbf{Trapezoidal error bound.} If $\abs{f''(x)} \le M_2$ for all $x$ in $[a,b]$, then
\[
\abs{\,\int_a^b f\,\dd x - T_n\,} \le \frac{(b-a)\,h^2}{12}\,M_2 = \frac{(b-a)^3}{12\,n^2}\,M_2.
\]
The error is controlled by the largest curvature on the interval and falls like $1/n^2$.
\end{focusbox}

\begin{labbox}
Bound the error of our $T_4$ estimate of $\int_0^1 e^{-x^2}\,\dd x$. We need a bound $M_2$ on $\abs{f''}$ over $[0,1]$. With $f = e^{-x^2}$, $f' = -2x e^{-x^2}$ and $f'' = (4x^2 - 2)e^{-x^2}$. On $[0,1]$ this ranges between $-2$ (at $x=0$) and about $0.74$ (at $x=1$), so $\abs{f''} \le 2$; take $M_2 = 2$. With $b-a = 1$, $n=4$, $h = 0.25$:
\[
\abs{\text{error}} \le \frac{1 \cdot (0.25)^2}{12}\cdot 2 = \frac{0.0625}{12}\cdot 2 = 0.0104.
\]
Our actual error was $0.0038$, comfortably inside the guaranteed bound of $0.0104$. The bound is valid but conservative -- it uses the worst-case curvature everywhere, while the real curvature is smaller over most of the interval. This is the typical relationship: the bound certifies, the true error beats it.
\end{labbox}

\section{Simpson's Rule: Let the Tops Be Parabolas}
\label{ch27:simpson}

The trapezoidal rule replaced the curve by straight chords and was left with an error driven by curvature, $f''$. The obvious next move: replace the curve not by lines but by \emph{parabolas}, which can themselves curve. A parabola is determined by three points, so we take the samples two slices at a time -- $x_0, x_1, x_2$ -- fit the unique parabola through those three heights, and use the exact area under that parabola as our estimate for the area over $[x_0, x_2]$ (Figure~\ref{ch27:fig-simp}).

\begin{figure}[H]
\centering
\begin{tikzpicture}
\begin{axis}[
  width=11.5cm, height=6cm,
  domain=0:4, samples=100,
  axis lines=middle,
  xtick={0,2,4}, xticklabels={$x_0$,$x_1$,$x_2$}, ytick=\empty,
  xlabel={$x$}, ylabel={$f(x)$},
  xmin=-0.3, xmax=4.6, ymin=0, ymax=3.4,
  enlargelimits=false,
]
\addplot[curveblue, very thick, domain=0:4, samples=80] {0.7 + 0.9*sin(deg(0.55*x)) + 0.18*x};
\addplot[targetgreen, very thick, dashed, domain=0:4, samples=60]
  {0.7 + (1.862-0.7)*(x/2) + ((2.148 - 2*1.862 + 0.7)/2)*(x/2)*(x/2-1)};
\addplot[only marks, mark=*, mark size=2pt, trapred] coordinates {(0,0.7) (2,1.862) (4,2.148)};
\node[curveblue, font=\small, anchor=south west] at (axis cs:0.35, 1.35) {curve $f(x)$};
\node[targetgreen!80!black, font=\small, anchor=south east] at (axis cs:3.9, 2.5) {parabola $p(x)$};
\end{axis}
\end{tikzpicture}
\caption{Simpson's rule fits a parabola (dashed green) through three consecutive samples and integrates the parabola exactly. A parabola can bend, so it tracks a curving function far better than a straight chord. The fit need not pass under or over the curve everywhere -- it weaves through it, and the over- and under-estimates largely cancel.}
\label{ch27:fig-simp}
\end{figure}

The arithmetic of ``integrate the parabola through three points'' produces a startlingly tidy weight pattern. Over a double-slice $[x_0, x_2]$ of total width $2h$, the exact area under the interpolating parabola (the unique quadratic passing through the three points) works out to
\[
\int_{x_0}^{x_2} (\text{parabola})\,\dd x = \frac{h}{3}\Big[\,f(x_0) + 4 f(x_1) + f(x_2)\,\Big].
\]
The middle point gets weight $4$, the two ends weight $1$, scaled by $h/3$. (The derivation is in the box below -- it is just integrating a quadratic.) Now stitch these double-slices together along $[a,b]$. This requires $n$ to be \emph{even} so the slices pair up cleanly. Each interior odd-indexed node $x_1, x_3, x_5, \dots$ is the apex of its own parabola and gets weight $4$. Each interior even-indexed node $x_2, x_4, \dots$ is shared as the right end of one parabola and the left end of the next, so its two weight-$1$ contributions add to $2$. The outer ends keep weight $1$.

\begin{focusbox}
\textbf{Simpson's Rule} (requires $n$ even). With $h = (b-a)/n$,
\[
S_n = \frac{h}{3}\Big[\, f(x_0) + 4f(x_1) + 2f(x_2) + 4f(x_3) + 2f(x_4) + \cdots + 4f(x_{n-1}) + f(x_n)\,\Big].
\]
The weight pattern is $1,4,2,4,2,\dots,4,1$: ends are $1$, odd interior nodes are $4$, even interior nodes are $2$, all scaled by $h/3$.
\end{focusbox}

\begin{whybox}
\textbf{Deriving the $\tfrac{h}{3}[1,4,1]$ formula.} Center the double-slice at the origin: nodes at $-h, 0, +h$ with heights $y_{-1}, y_0, y_1$. Any parabola is $p(x) = A + Bx + Cx^2$. Integrate it across $[-h,h]$:
\[
\int_{-h}^{h} (A + Bx + Cx^2)\,\dd x = 2Ah + \frac{2Ch^3}{3},
\]
since the odd term $Bx$ integrates to zero. Now express $A$ and $C$ through the sampled heights. Plugging in the nodes: $y_0 = A$, and $y_1 + y_{-1} = 2A + 2Ch^2$, so $C = (y_1 - 2y_0 + y_{-1})/(2h^2)$. Substitute:
\[
2Ah + \frac{2Ch^3}{3} = 2y_0 h + \frac{h}{3}(y_1 - 2y_0 + y_{-1}) = \frac{h}{3}\big(y_{-1} + 4y_0 + y_1\big).
\]
This yields the $1,4,1$ coefficient pattern. Notice that the linear coefficient $B$
does not appear in the integral -- the odd term integrates to zero by symmetry.
This symmetry cancellation is the primary reason for the high accuracy of Simpson's rule.
\end{whybox}

\begin{labbox}
Estimate $\int_0^1 e^{-x^2}\,\dd x$ with Simpson's rule, $n = 4$, reusing the exact same five heights from the trapezoidal lab. Apply the weights $1,4,2,4,1$:
\[
S_4 = \frac{0.25}{3}\Big[1.000000 + 4(0.939413) + 2(0.778801) + 4(0.569783) + 0.367879\Big].
\]
The bracket: $1.000000 + 3.757652 + 1.557602 + 2.279132 + 0.367879 = 8.962265$. Multiply by $h/3 = 0.083333$: $S_4 = 0.746855$. The true value is $0.746824$. The error is $0.000031$ -- about $31$ \emph{millionths}. Trapezoid with the very same data missed by $0.0038$. Simpson, using not a single extra function evaluation, is more than a hundred times more accurate. Let us see why that simple change in weights buys so much accuracy.
\end{labbox}

\section{Cubic Exactness and Fourth-Order Convergence}
\label{ch27:whysimpson}

Simpson's rule is derived by fitting parabolas, which guarantees exact integration
for quadratic polynomials. However, it is also mathematically exact for all \emph{cubic}
polynomials. That additional degree of polynomial exactness accounts for its fourth-order
convergence rate.

\begin{intuitionbox}
\textbf{Cubic exactness via symmetry.} Consider $f(x) = x^3$ on the symmetric double-slice $[-h, h]$ with midpoint at $0$. Its true integral is zero, because $x^3$ is an odd function: the positive area on $[0,h]$ exactly cancels the negative area on $[-h,0]$. Now evaluate Simpson's rule on this interval: the weights $1,4,1$ act on the sample values $(-h)^3 = -h^3$, $0^3 = 0$, and $h^3$, giving $1\cdot(-h^3) + 4\cdot 0 + 1\cdot h^3 = 0$. Simpson's rule yields zero with zero error. The symmetric evaluation cancels the cubic contribution just as the true integral does. Consequently, a method constructed to interpolate quadratics integrates all cubics exactly as well.
\end{intuitionbox}

\begin{recoverybox}
\textbf{Pause and think:} Suppose you apply Simpson's rule to $f(x) = 2x^3 - 5x^2 + 4x - 1$ on $[-2, 2]$ using $n = 2$. How far will your numerical estimate be from the exact integral?

\emph{Mental check:} The error is exactly zero. Simpson's rule fits a degree-2 parabola, but because the weighting is symmetric, the cubic and linear terms cancel out completely. Any polynomial up to degree 3 is integrated with zero error by Simpson's rule.
\end{recoverybox}

To understand this error cancellation, we examine the \emph{midpoint rule}. Instead of sampling at the ends of a slice (trapezoid) it samples at the \emph{center} and uses a rectangle of that single height. Evaluating at the midpoint changes the error behavior significantly: because the rectangle's flat top crosses the curve in the middle of the slice, it overshoots on one side and undershoots on the other, and those two errors largely balance each other (Figure~\ref{ch27:fig-mid}). That cancellation is why the midpoint rule, despite using only \emph{one} sample per slice, achieves $h^2$ accuracy -- and its leading error term has the \emph{opposite sign} of the trapezoidal error, which Simpson's rule exploits directly.

\begin{figure}[H]
\centering
\begin{tikzpicture}
\begin{axis}[
  width=11.5cm, height=6cm,
  domain=0:4, samples=100,
  axis lines=middle,
  xtick={0,1,2,3,4}, ytick=\empty,
  xlabel={$x$}, ylabel={$f(x)$},
  xmin=-0.3, xmax=4.5, ymin=0, ymax=3.2,
  enlargelimits=false,
]
\foreach \i in {0,1,2,3} {
  \edef\temp{\noexpand\addplot[draw=structurepurple!70, fill=structurepurple!12]
    coordinates {(\i,0) (\i+1,0) (\i+1,{0.6+0.25*(\i+0.5)+0.12*(\i+0.5)*(\i+0.5)}) (\i,{0.6+0.25*(\i+0.5)+0.12*(\i+0.5)*(\i+0.5)})} --cycle;}
  \temp
}
\addplot[curveblue, very thick, domain=0:4] {0.6 + 0.25*x + 0.12*x^2};
\addplot[only marks, mark=*, mark size=1.8pt, trapred] coordinates {
  (0.5,{0.6+0.25*0.5+0.12*0.25}) (1.5,{0.6+0.25*1.5+0.12*2.25})
  (2.5,{0.6+0.25*2.5+0.12*6.25}) (3.5,{0.6+0.25*3.5+0.12*12.25})};
\node[curveblue, font=\small, anchor=south east] at (axis cs:2.9, 2.45) {$f(x)$};
\end{axis}
\end{tikzpicture}
\caption{The midpoint rule: each rectangle's height is sampled at the \emph{center} of its slice (red dots). The flat top now slices through the curve, overshooting on the low side and undershooting on the high side, so the two errors largely cancel. One sample per slice, yet accuracy $\sim h^2$ -- and its leftover error tilts opposite to the trapezoid's, the fact Simpson exploits.}
\label{ch27:fig-mid}
\end{figure}

\begin{whybox}
\textbf{Simpson is the trapezoid's error, cancelled.} Recall the two facts we proved earlier on a single slice: the \emph{midpoint rule} (rectangle of height $f$ at the center) has error $+\tfrac{1}{24} f'' h^3$ per slice, while the \emph{trapezoid} has error $-\tfrac{1}{12}f'' h^3$ -- the trapezoid's leading error is exactly $-2$ times the midpoint's. So if we take a weighted blend that kills the $f''$ term,
\[
\frac{2\,(\text{Midpoint}) + 1\,(\text{Trapezoid})}{3},
\]
the leading curvature errors cancel: $2\cdot(+\tfrac{1}{24}) + 1\cdot(-\tfrac{1}{12}) = \tfrac{1}{12} - \tfrac{1}{12} = 0$. And this blend \emph{is} Simpson's rule. On a double-slice the trapezoid uses the two endpoints and the midpoint rule uses the center; combine them in the ratio $2{:}1$ and you reproduce exactly the $1,4,1$ pattern. So Simpson is not an unrelated trick: it is precisely the linear combination of the two simple rules engineered to annihilate their common $f''$ error term.

With the curvature error annihilated, where does the leading mistake go? Expanding $f$ to fourth order about the center of $[-h, h]$, the true integral is $2f(0)h + \tfrac{1}{3}f''(0)h^3 + \tfrac{1}{60}f^{(4)}(0)h^5 + \cdots$, while Simpson's formula evaluates to $2f(0)h + \tfrac{1}{3}f''(0)h^3 + \tfrac{1}{36}f^{(4)}(0)h^5 + \cdots$. The difference is $(\tfrac{1}{60} - \tfrac{1}{36})f^{(4)}(0)h^5 = -\tfrac{1}{90}f^{(4)}(0)h^5$ per double-slice of width $2h$. Summing over all $n/2$ double-slices yields a global error of $-\tfrac{(b-a)h^4}{180}f^{(4)}(\xi)$. Every lower order of error has completely vanished!
\end{whybox}

Because the $f''$ error is gone, Simpson's per-slice error is $\sim h^5 f^{(4)}$, and summed over the $n/2$ double-slices the global error is $\sim h^4 f^{(4)}$. Compare:

\begin{center}
\renewcommand{\arraystretch}{1.3}
\begin{tabular}{lccc}
\toprule
Rule & Exact for & Global error & Halving $h$ divides error by\\
\midrule
Rectangles & constants (deg 0) & $\sim h$ & $2$\\
Trapezoid  & lines (deg 1)     & $\sim h^2$ & $4$\\
Simpson    & cubics (deg 3)    & $\sim h^4$ & $16$\\
\bottomrule
\end{tabular}
\end{center}

These three growth laws are not abstractions; you can \emph{watch} them. Figure~\ref{ch27:fig-conv} plots the true error of each rule against the number of slices $n$ for our running integral $\int_0^1 e^{-x^2}\,\dd x$, on log--log axes (where each major grid line marks a power of ten in slices or error). On such axes a power law $\text{error} \sim n^{-p}$ becomes a straight line of slope $-p$, and the three rules sort themselves into three visibly different slopes.

\begin{figure}[H]
\centering
\begin{tikzpicture}
\begin{loglogaxis}[
  width=10.5cm, height=7.2cm,
  xlabel={number of slices $n$}, ylabel={absolute error},
  xmin=1.6, xmax=160, ymin=1e-10, ymax=1,
  grid=both, grid style={gray!18},
  legend pos=outer north east,
  legend cell align=left,
  log basis x=10, log basis y=10,
]
\addplot[trapred, very thick, mark=*, mark size=1.6pt] coordinates {
  (2,1.30e-1) (4,6.50e-2) (8,3.25e-2) (16,1.62e-2) (32,8.12e-3) (64,4.06e-3) (128,2.03e-3)
};
\addlegendentry{rectangles ($\sim n^{-1}$)}
\addplot[curveblue, very thick, mark=square*, mark size=1.6pt] coordinates {
  (2,1.55e-2) (4,3.84e-3) (8,9.59e-4) (16,2.40e-4) (32,5.99e-5) (64,1.50e-5) (128,3.74e-6)
};
\addlegendentry{trapezoid ($\sim n^{-2}$)}
\addplot[targetgreen, very thick, mark=triangle*, mark size=2pt] coordinates {
  (2,3.05e-4) (4,3.10e-5) (8,2.02e-6) (16,1.28e-7) (32,8.04e-9) (64,5.03e-10)
};
\addlegendentry{Simpson ($\sim n^{-4}$)}
\end{loglogaxis}
\end{tikzpicture}
\caption{Error of each rule versus the number of slices $n$ for $\int_0^1 e^{-x^2}\,\dd x$, on log--log axes, where a power law appears as a straight line of slope $-p$. Rectangles fall one decade per decade ($p=1$), the trapezoid two ($p=2$), Simpson four ($p=4$). The vertical gap is the whole story of the chapter: at $n=32$, Simpson is already near the limits of single precision (roughly seven decimal digits of precision) while the trapezoid still carries five-digit error and rectangles barely three.}
\label{ch27:fig-conv}
\end{figure}

\begin{focusbox}
\textbf{Simpson error bound.} If $\abs{f^{(4)}(x)} \le M_4$ on $[a,b]$, then
\[
\abs{\,\int_a^b f\,\dd x - S_n\,} \le \frac{(b-a)\,h^4}{180}\,M_4 = \frac{(b-a)^5}{180\,n^4}\,M_4.
\]
The error falls like $1/n^4$. Doubling the number of slices divides the error by $16$.
\end{focusbox}

\begin{labbox}
Bound the error of $S_4$ for $\int_0^1 e^{-x^2}\,\dd x$. We need $M_4 = \max\abs{f^{(4)}}$ on $[0,1]$. For $f = e^{-x^2}$, $f^{(4)} = (16x^4 - 48x^2 + 12)e^{-x^2}$. At $x = 0$ this is $12$; it decreases over $[0,1]$, so $\abs{f^{(4)}} \le 12$; take $M_4 = 12$. With $b-a=1$, $h=0.25$:
\[
\abs{\text{error}} \le \frac{1\cdot (0.25)^4}{180}\cdot 12 = \frac{0.00390625}{180}\cdot 12 = 0.00026.
\]
Our actual error was $0.000031$, well inside the certified $0.00026$. Even the conservative bound promises four good digits from five function values; the truth delivers five.
\end{labbox}

\begin{aestheticbox}
Moving from straight chords to parabolas -- with symmetry cancelling the cubic term -- improves the global error from $h^2$ to $h^4$. In practice, that difference in scaling is dramatic: to achieve an accuracy of $10^{-8}$, the trapezoidal rule requires on the order of $10^4$ subintervals, whereas Simpson's rule requires only on the order of $10^2$. Using identical function evaluations, adjusting the quadrature weights from $1,2,2,\dots,2,1$ to $1,4,2,4,\dots,4,1$ increases accuracy by orders of magnitude. In scientific computing, choosing a higher-order approximation method generally yields greater gains than simply increasing the number of grid points.
\end{aestheticbox}

\section{A Fuller Worked Example with Both Rules}
\label{ch27:worked}

Let us integrate something with a known answer so we can check ourselves against the truth. Take
\[
\int_1^2 \frac{1}{x}\,\dd x = \ln 2 = 0.693147\ldots
\]
This one \emph{does} have an antiderivative, $\ln x$, but that makes it a perfect test bench: we can compare our estimates against the exact value to the last digit.

Use $n = 4$, so $h = 0.25$ and the nodes are $1, 1.25, 1.5, 1.75, 2$. The heights $f(x) = 1/x$:
\[
\begin{array}{c|c}
x & 1/x\\\hline
1.00 & 1.000000\\
1.25 & 0.800000\\
1.50 & 0.666667\\
1.75 & 0.571429\\
2.00 & 0.500000
\end{array}
\]

\textbf{Trapezoidal.} Weights $1,2,2,2,1$:
\[
T_4 = \frac{0.25}{2}\big[1.000000 + 2(0.800000) + 2(0.666667) + 2(0.571429) + 0.500000\big].
\]
Bracket: $1.000000 + 1.600000 + 1.333334 + 1.142858 + 0.500000 = 5.576192$. Times $0.125$: $T_4 = 0.697024$. Error: $0.003877$.

\textbf{Simpson.} Weights $1,4,2,4,1$:
\[
S_4 = \frac{0.25}{3}\big[1.000000 + 4(0.800000) + 2(0.666667) + 4(0.571429) + 0.500000\big].
\]
Bracket: $1.000000 + 3.200000 + 1.333334 + 2.285716 + 0.500000 = 8.319050$. Times $0.083333$: $S_4 = 0.693254$. Error: $0.000107$.

Simpson beats the trapezoidal rule by a factor of about $36$ here, on identical data. Now let us check the error bounds. For $f = 1/x$ on $[1,2]$: $f'' = 2/x^3$, maximized at $x=1$ giving $M_2 = 2$; and $f^{(4)} = 24/x^5$, maximized at $x=1$ giving $M_4 = 24$.
\[
\text{Trap bound} = \frac{1\cdot(0.25)^2}{12}\cdot 2 = 0.01042,\qquad
\text{Simpson bound} = \frac{1\cdot(0.25)^4}{180}\cdot 24 = 0.00052.
\]
Both actual errors ($0.0039$ and $0.0001$) sit safely below their bounds. The numbers tell a consistent story across two completely different integrands.

\begin{labbox}
\textbf{Use the bound to choose $n$ in advance.} Suppose you want $\int_1^2 \tfrac1x\,\dd x$ to within $10^{-6}$ using Simpson. Solve the bound for $n$:
\[
\frac{(b-a)^5}{180\,n^4}M_4 \le 10^{-6}
\;\Longrightarrow\;
\frac{1\cdot 24}{180\,n^4} \le 10^{-6}
\;\Longrightarrow\;
n^4 \ge \frac{24}{180\cdot 10^{-6}} = 133{,}333
\;\Longrightarrow\;
n \ge 19.1.
\]
Round up to the next even integer: $n = 20$. With just twenty-one function values, Simpson guarantees six correct digits of $\ln 2$. The same calculation for the trapezoidal rule demands $n \ge 409$ -- twenty times the work for the same accuracy. This is how production code actually decides its step size: not by trial and error, but by inverting the error bound.
\end{labbox}

\begin{recoverybox}
\textbf{Pause and think:} Suppose your simulation needs to improve an integral's accuracy by a factor of $10{,}000$. How many times more slices must the trapezoidal rule use? How many times more for Simpson?

\emph{Mental check:} The trapezoid error falls like $1/n^2$, so cutting error by $10{,}000$ requires $\sqrt{10{,}000} = 100$ times as many slices. Simpson's error falls like $1/n^4$, so it requires only $(10{,}000)^{1/4} = 10$ times as many slices. A ten-fold increase in evaluation points buys a ten-thousand-fold leap in precision!
\end{recoverybox}

\begin{trapbox}
Three common numerical integration pitfalls:
\begin{itemize}
\item \textbf{Odd $n$ in Simpson.} Simpson pairs slices into parabolas, so $n$ must be even. With odd $n$ the weight pattern $1,4,2,\dots$ does not close properly and your answer is silently wrong. Always check $n$ is even before applying $h/3\,[1,4,2,\dots,4,1]$.
\item \textbf{Forgetting the rules assume smoothness.} The error bounds rely on $f''$ (trapezoid) or $f^{(4)}$ (Simpson) being bounded. If $f$ has a kink, a spike, or a vertical tangent in the interval -- think $\sqrt{x}$ near $0$, whose derivatives blow up -- the bound is infinite and the promised convergence rate evaporates. Subdivide at the bad point or change variables first.
\item \textbf{Mixing up $h/2$ and $h/3$.} Trapezoid is $h/2$, Simpson is $h/3$. Swap them and every digit is off. The mnemonic: trapezoid \emph{averages} two heights ($\div 2$); Simpson's weights $1{+}4{+}1 = 6 = 3\times 2$ live naturally over $\div 3$.
\end{itemize}
\end{trapbox}

\section{How Computers Actually Integrate}
\label{ch27:computers}

When you type \texttt{quad(f, a, b)} in a scientific library, what runs underneath is a refined descendant of exactly what we have built. The core ideas survive; engineering wraps them.

\textbf{Adaptive subdivision.} A fixed, uniform $h$ is wasteful. Where $f$ is nearly straight, big slices suffice; where $f$ swings violently, you want tiny slices. So real integrators are \emph{adaptive}: they apply Simpson's rule on a region, then split the region in half and apply Simpson to each half. If the two-half estimate agrees with the whole-region estimate to within tolerance, the region is declared done; otherwise each half is recursively subdivided. Effort flows automatically to wherever the function is hardest, which is the only sensible place to spend it (Figure~\ref{ch27:fig-adaptive}).

\begin{figure}[H]
\centering
\begin{tikzpicture}
\begin{axis}[
  width=11.5cm, height=6cm,
  domain=0:4, samples=100,
  axis lines=middle,
  xtick={0, 1.5, 4.0},
  xticklabels={$a$, $x_1$, $b$},
  ytick=\empty,
  xlabel={$x$}, ylabel={$f(x)$},
  xmin=-0.2, xmax=4.5, ymin=0, ymax=2.8,
  enlargelimits=false,
]
\addplot[draw=targetgreen!70, fill=targetgreen!12] coordinates {
  (0,0) (1.5,0) (1.5,0.48) (0,0.40)
} --cycle;
\addplot[draw=targetgreen!70, fill=targetgreen!12] coordinates {
  (1.5,0) (2.2,0) (2.2,1.14) (1.5,0.48)
} --cycle;
\addplot[draw=targetgreen!70, fill=targetgreen!12] coordinates {
  (2.2,0) (2.5,0) (2.5,1.97) (2.2,1.14)
} --cycle;
\addplot[draw=targetgreen!70, fill=targetgreen!12] coordinates {
  (2.5,0) (2.7,0) (2.7,2.24) (2.5,1.97)
} --cycle;
\addplot[draw=targetgreen!70, fill=targetgreen!12] coordinates {
  (2.7,0) (2.9,0) (2.9,1.99) (2.7,2.24)
} --cycle;
\addplot[draw=targetgreen!70, fill=targetgreen!12] coordinates {
  (2.9,0) (3.2,0) (3.2,1.19) (2.9,1.99)
} --cycle;
\addplot[draw=targetgreen!70, fill=targetgreen!12] coordinates {
  (3.2,0) (4.0,0) (4.0,0.60) (3.2,1.19)
} --cycle;
\addplot[curveblue, very thick, domain=0:4, samples=120] {0.4 + 0.05*x + 1.7*exp(-4*(x-2.7)^2)};
\addplot[only marks, mark=*, mark size=1.8pt, trapred] coordinates {
  (0,0.40) (1.5,0.48) (2.2,1.14) (2.5,1.97) (2.7,2.24) (2.9,1.99) (3.2,1.19) (4.0,0.60)
};
\node[font=\footnotesize, fill=white, inner sep=1pt] at (axis cs:0.75, 1.2) {wide slices ($h$ large)};
\draw[->, thick, gray] (axis cs:0.75, 1.05) -- (axis cs:0.75, 0.65);
\node[font=\footnotesize, fill=white, inner sep=1pt] at (axis cs:2.7, 2.55) {dense slices ($h$ tiny)};
\end{axis}
\end{tikzpicture}
\caption{Adaptive quadrature in action. Where the curve is gentle and nearly straight (left), a wide interval suffices. Where the curve surges into a sharp peak (center-right), the algorithm recursively bisects the region into narrow slices to capture the high curvature. Computation is saved in the flats and concentrated at the peak.}
\label{ch27:fig-adaptive}
\end{figure}

\begin{whybox}
\textbf{How the computer estimates its own error without knowing the answer.} This is the cleverest practical idea in the chapter. Run Simpson once on $[a,b]$ with step $h$ to get $S(h)$, then again with step $h/2$ to get $S(h/2)$. We know the error scales like $h^4$, so $S(h)$ is off by about $C h^4$ and $S(h/2)$ by about $C(h/2)^4 = Ch^4/16$. Subtract the two estimates:
\[
S(h/2) - S(h) \approx \Big(I - \tfrac{Ch^4}{16}\Big) - \big(I - Ch^4\big) = Ch^4 - \frac{Ch^4}{16} = \frac{15}{16}\,Ch^4.
\]
The thing we cannot see directly -- the error $Ch^4/16$ of the better estimate -- is therefore about $\tfrac{1}{15}$ of the \emph{difference between the two estimates}, which we \emph{can} see. So the computer gauges its own error by comparing two grids, with no oracle for the true value. When that estimated error drops below your tolerance, it stops. This is the engine inside every adaptive quadrature routine.
\end{whybox}

\begin{connectionbox}
That same comparison-of-two-grids trick can be turned into \emph{Richardson extrapolation}: form the combination $\big(16\,S(h/2) - S(h)\big)/15$, which cancels the leading $h^4$ error entirely and leaves an error of order $h^6$. Applied systematically to the trapezoidal rule, this idea generates the \emph{Romberg integration} table -- a triangular scheme that raises the accuracy order at each column. Romberg's first column is trapezoidal estimates on ever-finer grids; its second column is exactly Simpson's rule; later columns are higher-order rules nobody had to derive by hand. It is the same Taylor-cancellation logic used to construct Simpson's rule, carried to higher orders. We meet this extrapolation principle again when studying Taylor series.
\end{connectionbox}

\textbf{Gaussian quadrature.} We assumed the sample points were \emph{equally spaced}, because that is natural and because it is forced on you when $f$ is a table of measurements. But if you are free to choose \emph{where} to sample $f$, you can do dramatically better. Gauss asked: with a fixed budget of $n$ sample points, placed wherever we like and weighted however we like, what is the most accurate rule possible? The answer -- sampling at special points called the roots of Legendre polynomials, clustered toward the ends of the interval -- yields a rule exact for polynomials up to degree $2n-1$. With the same $n$ evaluations that give Simpson degree-$3$ exactness, Gaussian quadrature reaches degree $2n-1$. For smooth functions this is the gold standard, and it is what high-precision libraries reach for first.

\begin{historybox}
Around 250~BCE, Archimedes approximated the area of a circle using inscribed and
circumscribed polygons, establishing the principle of bounding areas by simpler shapes.
Surveys in Egypt and Rome similarly estimated irregular land areas using trapezoidal slices.

The systematic theory began around 1676--1711, when Isaac Newton and Roger Cotes
developed what is now called \textbf{Newton--Cotes quadrature}. Working on planetary
orbits where no algebraic antiderivatives were available, they interpolated functions
with polynomials over equally spaced points.

Thomas Simpson published the three-point parabolic rule bearing his name in his 1743
treatise \emph{Mathematical Dissertations}, though the formula had been derived earlier
by James Gregory in 1668 and used by Newton and Cavalieri.

Newton and Cotes assumed that increasing polynomial degree on equally spaced nodes would
yield arbitrary accuracy. However, for degree 8 and higher, the interpolation weights
oscillate in sign and grow large, making the formulas sensitive to rounding errors
(a manifestation of Runge's phenomenon). In 1814, Carl Friedrich Gauss showed that by
choosing both the sample locations and their weights optimally, an $n$-point rule could
integrate polynomials up to degree $2n-1$ exactly, with all weights positive. Modern
scientific software commonly uses adaptive extensions such as Gauss--Kronrod quadrature.
\end{historybox}

\begin{happeningbox}
\textbf{What a computer evaluates} during numerical quadrature is a concrete weighted sum. It samples $f$ at a set of designated points, multiplies those function values by precomputed quadrature weights, and sums the results. Every definite integral evaluated computationally -- a probability, a Fourier coefficient, the energy of an orbit, or a machine learning loss -- reduces to a weighted sum of function values: $\sum_i w_i f(x_i)$. The trapezoidal weights are $1,2,2,\dots,2,1$; Simpson's are $1,4,2,4,\dots,1$; Gaussian quadrature weights are derived from orthogonal polynomials. In each case, numerical integration selects nodes $x_i$ and weights $w_i$ so that the finite sum approximates continuous accumulation as accurately as the smoothness of $f$ allows.
\end{happeningbox}

\begin{walkbox}
On your walk, picture the speedometer-and-distance problem. You glance at the speedometer once a minute on a winding drive. The trapezoidal rule says: average each consecutive pair of readings and multiply by a minute. Simpson says: take readings three at a time and trust that the speed varied parabolically between them. Now ask yourself: on a stretch where your acceleration was steadily increasing (constant $f''$ but $f''' = 0$), which rule is exactly right, and why does Simpson nail it while the trapezoid does not? Then ask the harder one: if you got to choose \emph{when} to glance, would you space your glances evenly -- or look more often through the twists and less often on the straightaways?
\end{walkbox}

\section{Summary: The Logic of the Whole Subject}
\label{ch27:summary}

We began stranded -- handed integrals like $\int_0^1 e^{-x^2}\,\dd x$ that provably have no elementary antiderivative, so the Fundamental Theorem cannot help. We returned to the integral's true meaning as area and learned to estimate that area from samples.

\begin{itemize}
\item \textbf{Rectangles} approximate $f$ by constants on each slice. Global error $\sim h$. Crude.
\item \textbf{Trapezoidal rule}, weights $\tfrac{h}{2}[1,2,\dots,2,1]$, approximates $f$ by straight chords. Exact for lines; error driven by curvature $f''$; global error $\sim h^2$. Built from connect-the-dots.
\item \textbf{Simpson's rule}, weights $\tfrac{h}{3}[1,4,2,4,\dots,4,1]$ with $n$ even, approximates $f$ by parabolas. Exact for cubics (one bonus degree, by symmetry); error driven by $f^{(4)}$; global error $\sim h^4$. Built as the curvature-cancelling blend $\tfrac{2\,\text{Mid} + \text{Trap}}{3}$.
\item \textbf{Error bounds} let you certify accuracy and choose $n$ in advance: trapezoid $\le \tfrac{(b-a)^3}{12n^2}M_2$, Simpson $\le \tfrac{(b-a)^5}{180n^4}M_4$.
\item \textbf{Real software} uses these as building blocks inside adaptive subdivision (spend effort where $f$ is wild), self-error-estimation by comparing grids, Richardson extrapolation and Romberg tables, and Gaussian quadrature (optimal sample placement). It is always, in the end, a weighted sum $\sum w_i f(x_i)$.
\end{itemize}

The throughline is a single idea, applied with increasing refinement: \emph{replace the curve, locally, by something you can integrate exactly, and the better the local fit, the faster the error dies}. Constant, line, parabola -- and the error exponent climbs $h, h^2, h^4$. That is calculus turning its own definition into an algorithm.

\section{Limitations and the road ahead}
\label{ch27:reach}

\begin{connectionbox}
Numerical quadrature evaluates definite integrals directly from samples, bypassing
the need for elementary antiderivatives. When formulas like Simpson's rule or
Gauss--Kronrod quadrature are applied to smooth functions in one dimension, accurate
approximations require only a modest number of function evaluations.

However, standard grid-based quadrature faces a major limitation in higher dimensions.
When integrating over a $d$-dimensional domain, a grid of $N$ points per axis requires
$N^d$ total function evaluations:
\begin{itemize}
  \item In $1$ dimension: $100^1 = 100$ evaluations;
  \item In $3$ dimensions: $100^3 = 10^6$ evaluations;
  \item In $10$ dimensions: $100^{10} = 10^{20}$ evaluations.
\end{itemize}
This exponential growth in computation is known as the \emph{curse of dimensionality}.
Deterministic grids quickly become impractical as dimension grows.

In high-dimensional problems, two alternative approaches are standard:
\begin{itemize}
  \item \emph{Analytic reductions:} When symmetry exists, coordinate transformations
  and Fubini's theorem can reduce multidimensional integrals to simpler nested forms
  (Chapter~34).
  \item \emph{Monte Carlo integration:} In high dimensions, random sampling estimates
  integrals with an error of $O(1/\sqrt{N})$ independent of the dimension $d$.
\end{itemize}
\end{connectionbox}

\subsection{From Continuous Integrals to Discrete Lists: The Bridge to Chapter~28}
\label{ch27:bridge}

Evaluating $\int_0^1 e^{-x^2}\,\dd x$ by repeatedly halving the step size $h = 1/n$
generated a succession of trapezoidal estimates:
\[
  T_1,\; T_2,\; T_4,\; T_8,\; T_{16},\; T_{32},\; \dots,\; T_{2^k},\; \dots
\]
Writing out their numerical values reveals an orderly approach toward the true area:
\[
  0.683940,\quad 0.731370,\quad 0.742984,\quad 0.745866,\quad 0.746584,\quad 0.746764,\quad \dots
\]
This highlights a conceptual transition. The first twenty-seven chapters focused on the
continuum, where variables vary continuously. Numerical approximations generate an ordered
list of discrete numbers:
\[
  a_1,\; a_2,\; a_3,\; a_4,\; \dots,\; a_k,\; \dots
\]
where each entry $a_k = T_{2^k}$ is indexed by a discrete integer $k \in \N$. As the grid
is refined, this list approaches a definite limit:
\[
  \lim_{k \to \infty} a_k = \int_0^1 e^{-x^2}\,\dd x.
\]
Such an infinite ordered list is a \emph{sequence}.

Having studied continuous rates of change and continuous accumulation, we now turn to
discrete lists and infinite sums. In Chapter~28, we begin with \emph{Sequences and Their Limits}.

\section{Exercises}
\label{ch27:exercises}

\begin{exercise}[Mental checks]
Answer each of the following without pencil or paper:
\begin{enumerate}[label=(\alph*)]
\item If $f(x)$ is strictly concave down ($f'' < 0$) on $[a,b]$, will the trapezoidal estimate $T_n$ be strictly greater than or strictly less than the true integral $\int_a^b f(x)\,\dd x$?
\item You evaluate an integral using the trapezoidal rule with $n = 10$ and obtain an error of $0.04$. If you increase to $n = 20$, approximately what error do you expect?
\item You evaluate the same integral with Simpson's rule and find an error of $0.0016$ with $n = 10$. If you increase to $n = 20$, approximately what error do you expect?
\item For which of the following functions is Simpson's rule with $n = 2$ mathematically exact: $f(x) = 4x^3 - 2x + 7$, $g(x) = x^4$, $h(x) = \sin x$?
\end{enumerate}
\end{exercise}

\begin{exercise}
Estimate $\int_0^2 \sqrt{1 + x^3}\,\dd x$ using the trapezoidal rule with $n = 4$. (Nodes $0, 0.5, 1, 1.5, 2$; compute each $\sqrt{1+x^3}$ to four decimals, then apply weights $1,2,2,2,1$ and multiply by $h/2$.)
\end{exercise}

\begin{exercise}
Estimate the same integral $\int_0^2 \sqrt{1+x^3}\,\dd x$ with Simpson's rule, $n = 4$, reusing your five heights from Exercise 1 but with weights $1,4,2,4,1$ and factor $h/3$. By how much do the two estimates differ?
\end{exercise}

\begin{exercise}
Compute $\int_0^{\pi} \sin x\,\dd x$ exactly (you know the antiderivative), then estimate it with Simpson's rule using $n = 2$ (nodes $0, \pi/2, \pi$). Find the actual error. Why is it not zero even though $\sin$ is smooth?
\end{exercise}

\begin{exercise}
For $\int_0^1 e^{-x^2}\,\dd x$, how large must $n$ be to guarantee the trapezoidal rule is within $10^{-4}$ of the truth? Use $M_2 = 2$. Then find the $n$ needed for Simpson to reach the same tolerance using $M_4 = 12$. Comment on the ratio.
\end{exercise}

\begin{exercise}
A car's speed (in m/s) is recorded every $2$ seconds: at $t = 0,2,4,6,8,10,12$ the readings are $0, 8, 17, 22, 24, 25, 25$. This gives $n = 6$ intervals -- an even number, so Simpson applies. Estimate the distance travelled, $\int_0^{12} v\,\dd t$, with (a) the trapezoidal rule and (b) Simpson's rule. Which would you trust more, and why?
\end{exercise}

\begin{exercise}
(Conceptual.) Explain in your own words why the trapezoidal rule \emph{overestimates} the integral of a function that is concave up (like $e^x$) on the whole interval, and underestimates a concave-down function. Tie your answer to the sign of $f''$ in the error term $-\tfrac{h^3}{12}f''(\xi)$ per slice.
\end{exercise}

\begin{exercise}
(Proof.) Show directly that Simpson's rule is exact for every cubic $f(x) = ax^3 + bx^2 + cx + d$ on a single double-slice $[-h, h]$. That is, compute $\int_{-h}^{h} f\,\dd x$ exactly and compare it to $\tfrac{h}{3}[f(-h) + 4f(0) + f(h)]$, and verify they are equal. Identify which term of $f$ would have to be present (degree $4$) for Simpson to start being wrong.
\end{exercise}

\begin{exercise}
(Proof / derivation.) Show that for the trapezoidal rule applied to a single slice $[x_0, x_1]$ of width $h$, the leading error term is $-\tfrac{h^3}{12}f''(\xi)$ by Taylor-expanding $f$ about the midpoint and integrating, mirroring the symmetric derivation in the chapter.
\end{exercise}

\begin{exercise}
(Self-error estimation.) Using $\int_1^2 \tfrac1x\,\dd x$ with Simpson, compute $S_2$ (i.e. $h = 0.5$) and $S_4$ ($h = 0.25$). Form the estimated error of $S_4$ as $\tfrac{1}{15}(S_4 - S_2)$ and compare it to the actual error $\ln 2 - S_4$. How close is the estimate? This is exactly what adaptive software does.
\end{exercise}

\begin{exercise}
(Richardson extrapolation.) Continue the previous exercise: form the combination $S^\star = \tfrac{16 S_4 - S_2}{15}$, which cancels the leading $h^4$ error. Compute $S^\star$ and its error against $\ln 2 = 0.693147$. You should find it dramatically more accurate than either $S_2$ or $S_4$ alone -- this is the first step of Romberg integration.
\end{exercise}

\begin{exercise}[The sequence of approximations]
Let $f(x) = e^{-x^2}$ on $[0,1]$. For each positive integer $k \in \N$, let $s_k = S_{2^k}$ denote the Simpson's rule approximation with $n = 2^k$ slices.
\begin{enumerate}[label=(\alph*)]
\item Write down the first three terms of this list: $s_1 = S_2$, $s_2 = S_4$, $s_3 = S_8$.
\item Explain why the collection $(s_k)_{k=1}^\infty$ is a \emph{sequence} of real numbers, and what it means to write $\lim_{k \to \infty} s_k = \int_0^1 e^{-x^2}\,\dd x$.
\item Using the error bound $\abs{\int_0^1 e^{-x^2}\,\dd x - S_n} \le \frac{M_4}{180 n^4}$ with $M_4 = 12$, find a power $n = 2^k$ that mathematically guarantees $\abs{\int_0^1 e^{-x^2}\,\dd x - s_k} < 10^{-10}$. (This challenge-and-response behavior on discrete indices is the foundation of Chapter~28.)
\end{enumerate}
\end{exercise}
